problem:  
Let \( E_{p,q} = SU(3)//S^1_{p,q} \) be an Eschenburg space defined via the free \( S^1 \)-action  
\[
z \cdot A = \operatorname{diag}(z^{p_1}, z^{p_2}, z^{p_3})\, A \,\operatorname{diag}(\overline{z}^{\,q_1}, \overline{z}^{\,q_2}, \overline{z}^{\,q_3}),
\]
for integer vectors \( p, q \in \mathbb{Z}^3 \) satisfying \( \sum_i p_i = \sum_i q_i \).  
Certain such biquotients admit metrics of positive sectional curvature, while others only support nonnegative curvature.  

**Research problem (two facets):**  
1. **Existence question:**  
   Suppose that \( E_{p,q} \) admits a smooth, effective, isometric action by a 3-torus \( T^3 \).  
   Does the presence of such a high-rank torus symmetry imply the existence of some Riemannian metric on \( E_{p,q} \) with *strictly positive sectional curvature*?  

2. **Deformation question:**  
   If \( E_{p,q} \) admits a \( T^3 \)-invariant metric of *nonnegative* sectional curvature, can this metric be continuously deformed (within the space of \( T^3 \)-invariant metrics) to a metric of *positive* curvature?  

Why is it a "good" problem:  
This problem explicitly isolates two deep and complementary aspects — **existence** and **deformation** — concerning the interplay of symmetry and curvature on Eschenburg spaces. The first part probes whether large torus symmetry is itself a geometric rigidity condition enforcing positive curvature. The second part investigates the space of invariant metrics, testing whether positivity can be reached continuously from nonnegativity in a highly symmetric setting.  

Eschenburg spaces form an ideal class for such study, sitting at the interface of **homogeneous geometry and manifold topology**, with rich but tractable group actions. The problem thus extends the spirit of the **Grove–Searle symmetry-rank and positive curvature program** into the nonhomogeneous biquotient realm.  

The statement is conceptually simple—"Does enough symmetry force positivity?"—yet pursuing it demands sophisticated tools from **Lie group theory, Riemannian submersion geometry, metric deformation techniques, and global curvature analysis**. Its importance lies in potentially revealing new rigidity or flexibility phenomena for metrics on exotic 7‑manifolds, embodying the ideal of a “narrow entrance, profound depth” research problem.